\section{Glossary}\label{sec:glossary}

This appendix defines the domain terms as they are used throughout this whitepaper.
Terms appear in alphabetical order. Where a term names a concrete construct, the
relevant source location is cited against the current working tree; where a term names
a type or routine introduced by this design, it is identified as new and unimplemented.

\begin{description}

\item[Abut] A \seat{Abut} \code{SeatKind}: a rim-to-rim seat between two parts of
  equal radius (a tube fore rim meeting a tube aft rim). The radial check requires the
  two stations' outer radii to match within tolerance, and the gap is interpreted as a
  forward standoff along \zfwd, flush when the gap is zero. It is the default seat, so
  the trivial stack costs no authored numbers.

\item[CM-to-CM legacy link] The current attachment representation: each child stores a
  \code{Vector3} giving the child's center of mass relative to the parent's center of
  mass, in \src{model/parts/Part.h}{326}. Because a part's CM is not in general at its
  mid-length (a solid cone's CM sits $L/4$ from the base), this couples placement to a
  derived quantity and forces authors to hand-compute offsets. This design replaces it
  with a geometric \code{StationLink}.

\item[Composite CM / composite inertia] The center of mass and inertia tensor of a part
  together with all its descendants. They are cached on each part (\code{compositeCm} at
  \src{model/parts/Part.h}{320}, \code{compositeInertiaTensor} at
  \src{model/parts/Part.h}{300}) and recomputed by the
  two-pass routine \code{computeCompositeAt} (\src{model/parts/Part.cpp}{169}): pass one
  builds the mass-weighted composite CM, pass two shifts every child tensor to that CM via
  the parallel-axis theorem. The composite tensor is full mass-weighted ($\unit{kg\cdot m^2}$),
  whereas a part's own \code{inertiaTensor} is stored per unit mass ($\unit{m^2}$).

\item[Cone reversed-frame defect] A pre-existing implementation mistake in
  \code{coneCmOffset} (\src{model/parts/ConicalNoseCone.cpp}{49}): it reports the cone CM in a
  mid-length-origin, base-at-$+z$ frame, whose $+z$ is the \emph{opposite} of the chosen forward
  direction. This design \emph{corrects} it so the cone, like every part, reports its CM in the
  shared fore-plane-origin, \zfwd\ frame ($L/2 - \bar h \rightarrow \bar h - L/2$). The centroidal
  tensor is unchanged (axisymmetric, $z$-flip invariant); only the CM sign was wrong. With this,
  there is no frame adapter and no per-part CM override anywhere in the design.

\item[Envelope sweep] Layer~2 of the diagnostics (\code{sweepOverlaps}, new): after the
  whole tree resolves, build world axial intervals for every part, sort by aft station
  ($O(N\log N)$), then sample each offender at feature breakpoints and test its outer
  radius against each candidate host's capacity using the solid-host rule. It runs once
  per structural change, never per integration step.

\item[\code{getCenterMassOffset}] The accessor returning a part's CM relative to the
  component middle (\src{model/parts/Part.h}{96}), zero for symmetric parts and nonzero
  for a cone or fin set. It is documented but unconsumed today; this design consumes it
  through the uniform local-CM rule $-L/2 + \code{getCenterMassOffset}().z()$, the same for
  every part.

\item[\code{hostCapacity}] The radius a host occupies at a given station, defined by the
  solid-host rule as $\code{radiusInnerAt}(z)$ for a bored part and
  $\code{radiusOuterAt}(z)$ for a solid part. An offender of outer radius
  $r_{\mathrm{off}}$ fits iff $r_{\mathrm{off}} \le \code{hostCapacity}(z) + \mathit{tol}$.

\item[\code{innerCapacityAt}] The new non-virtual \code{Part} helper that evaluates
  \code{hostCapacity} for a part at a local station: the bore for a tube, the outer skin
  for a solid. One predicate, correct for both host kinds, used by the envelope sweep.

\item[\code{NestInBore}] A \seat{NestInBore} \code{SeatKind}: the child's outer radius
  seats inside the parent's bore. The radial check requires child OD within parent
  ID, and the gap is an insertion depth $\ge 0$ that drives the child \emph{aft} ($-z$)
  into the bore, so the signed gap is negated relative to \seat{Abut}.

\item[\code{OnSurface}] An \seat{OnSurface} \code{SeatKind}: the child seats radially on
  the parent's outer wall at a chosen station (a fin root on a body tube). The radial
  check compares the child's radius against the parent's outer radius at the station, and
  the gap is an axial standoff along the wall.

\item[\code{OverlapDiagnostic}] A new POD recording one located interference: the offender
  and host \code{Part::Id}s, the world station of the worst violation, the penetration
  depth in metres, and a human-readable message. \code{Part::Id} is
  \cppinline|std::uint64_t| (\src{model/parts/Part.h}{56}).

\item[Parallel-axis theorem (\code{parallelAxisTerm})] The displacement tensor
  $f(d) = (d\cdot d)\,I_3 - d\,d^{\mathsf T}$ that, scaled by a body's mass and added to
  its centroidal tensor, shifts that tensor to a parallel axis offset by $d$
  (\src{model/parts/Part.cpp}{17}). The composite inertia pass applies it to every child;
  this design preserves the math and changes only the input $d$ from CM-to-CM offsets to
  resolver-derived geometric stations.

\item[\code{Part} composite tree] The ownership tree of rocket components rooted at
  \code{Part} (\srcf{model/parts/Part.h}, namespace \code{model::part}). Each node is both
  a single component and the composite of itself with its children; children are stored in
  \code{childParts} (\src{model/parts/Part.h}{326}) and adopted by move via
  \code{addChildPart} (\src{model/parts/Part.h}{239}). \code{clone()} deep-copies and
  re-mints ids.

\item[\code{placementDirty} vs.\ the mass-delta gate] Two distinct cache gates. The
  existing mass-delta gate, \code{ensureCompositeCache} keyed on
  \code{builtAtCompositeMass} (\src{model/parts/Part.h}{288}), fires on every integration
  step during a burn because the composite mass changes each step. The new
  \code{placementDirty} flag would instead gate the resolver cache on \emph{structural}
  change only (a part added, removed, or geometrically edited). Today only the single
  \code{needsRecomputing} flag (\src{model/parts/Part.h}{322}) exists, propagated upward by
  \code{markAsNeedsRecomputing} (\src{model/parts/Part.h}{293}); the separate
  \code{placementDirty} and \code{resolvedCache} are added by this design.

\item[\code{Placed}] A new value pairing a resolved \code{Part} pointer with its
  \code{Pose}. The resolver returns a \cppinline|std::vector<Placed>| in deterministic
  depth-first attachment order.

\item[\code{Pose}] A new value type carrying an origin (the part-local fore-plane origin in
  the root frame) and an orientation quaternion (identity in 3-DOF). Its \code{compose}
  operator chains a child pose through its parent, rotating the child translation, so the
  same routine serves 6-DOF unchanged.

\item[\code{Propagatable}] The model$\leftrightarrow$sim bridge interface implemented by
  \code{RocketModel}, exposing \code{getForces}, \code{getTorques}, \code{getMass(t)}, and
  \code{terminateCondition} to the propagator. It is unaffected by this design, which
  changes only how composite quantities are derived, not how they are consumed by the
  integrator.

\item[Resolver (\code{resolvePlacements})] The new single authority for absolute
  placement: a depth-first walk over the ownership tree from a root pose, producing a
  \code{Placed} for every part by pure geometry (no CM, no time). Both the simulator
  (\code{computeCompositeAt}, \code{accumulateAeroAt}) and the visualizer
  (\code{RocketMesh::walk}) consume its output, so the two cannot disagree.

\item[SeatKind] A new closed enum of three radial seat relationships --- \seat{Abut},
  \seat{NestInBore}, \seat{OnSurface} --- that governs the radial compatibility check and
  the sign of the gap. It is intentionally closed because a rigid axisymmetric stack
  expresses only these three relationships.

\item[\code{SolveResult}] A new aggregate holding an \code{ok} flag and the list of
  \code{OverlapDiagnostic}s from a resolve. When \code{ok} is false, both the simulator and
  the visualizer refuse the solve; it is produced once per structural resolve and cached,
  so the gate costs nothing per integration step.

\item[Solid-host rule] The diagnostic predicate that resolves the ambiguity of a solid
  part, which has no bore and would otherwise flag everything. A host occupies
  $[0, \code{hostCapacity}(z)]$, where capacity is the bore for a bored part and the outer
  skin for a solid; an offender fits iff its outer radius does not exceed that capacity
  within tolerance. This makes the coupler-versus-cone test the correct poke-through check.

\item[Station landmark] The new \code{Station} value
  $\{z, r_{\mathrm{outer}}, r_{\mathrm{inner}}\}$ returned by \code{stationAt} for a
  fractional station: the live axial coordinate and radii in the part's \zfwd\ local frame,
  derived from geometry and never stored.

\item[Station / fractional station (\code{station01})] A position along a part's
  longitudinal axis expressed as a fraction in $[0,1]$, where $0$ is the aft plane and $1$
  the fore plane. The absolute axial coordinate is
  $(\code{station01}-1)\times\code{getLength()}$ (so $z \in [-L,0]$, with the aft plane at $-L$
  and the fore plane at the origin), recomputed on each read so that editing a
  length moves the seam. \code{getLength()} is promoted to a base \code{Part} virtual by
  this design; today it exists only on \code{ConicalNoseCone} and \code{BodyTube}.

\item[StationLink] The new stored attachment value replacing the CM-to-CM \code{Vector3}:
  a parent fractional station, a child fractional station, an axial gap, a \code{SeatKind},
  and a per-seam child rotation (identity in 3-DOF). It encodes physical intent --- which
  feature meets which --- rather than a derived coordinate, so the storage type becomes
  \cppinline|std::vector<std::pair<std::shared_ptr<Part>, StationLink>>|.

\item[3-DOF vs.\ 6-DOF] The simulator currently integrates only the three translational
  degrees of freedom; orientation is carried in \code{StateData} but not yet integrated.
  \code{Part} documents the rigid no-relative-rotation assumption
  (\src{model/parts/Part.h}{43}). This design shapes in the 6-DOF path at zero schema cost
  --- \code{Pose} carries an orientation, \code{StationLink} carries a per-seam rotation,
  and one composition line would rotate each child tensor --- while remaining bit-identical
  to today in 3-DOF.

\item[$+z = $ forward convention (\zfwd)] The settled longitudinal convention: $+z$ points
  toward the nose tip, every part's local origin sits on the axis at its \emph{fore} plane, and a
  part occupies $z \in [-\code{length}, 0]$. The fore origin makes the root's origin --- the
  nose tip --- the whole-assembly datum, so the derived CM is reported relative to the tip. The
  $+z$ \emph{direction} matches the visualizer and the existing \code{.qrd} files, so the renderer
  change is a consumer swap and the file migration a re-expression, both with no axial sign flip.
  \zaft\ was rejected for requiring exactly those inversions.

\end{description}
